\section{Introduction and Background}
\label{sec:introduction}

\subsection{Problem Statement}

Roughly one in seven adults in the United States meets criteria for an
alcohol-use disorder in any given year, yet only a small minority ever speak to
a counselor about it~\cite{niaaa2023}. The evidence-based clinical approach to
ambivalence about substance use is \textit{Motivational Interviewing} (MI),
developed by Miller and Rollnick~\cite{miller2013mi}. MI is collaborative and
goal-oriented but explicitly non-confrontational: rather than tell the patient
what to do, the clinician helps the patient articulate their own reasons for
change.

Two facts motivate this project. First, qualified MI clinicians do not scale --
training, supervision, and licensure are slow and expensive, and access is
geographically uneven. Second, off-the-shelf LLMs do not, by default, behave
like MI counselors. They are trained on instruction-following corpora that
reward direct answers and informative explanations, so when faced with a user
who says \textit{``I don't think I drink that much''}, the modal response from
a generic assistant is to either reassure the user (a missed opportunity to
develop discrepancy) or to enumerate health risks (a frontal assault that
provokes resistance). Neither is MI.

The question this capstone addresses is therefore practical:
\begin{quote}
\textit{Can a small, locally-deployable language model be turned into a
credible Motivational Interviewing assistant for alcohol-related conversations,
with explicit safety boundaries and full conversation auditability?}
\end{quote}

\subsection{Project Goals and Objectives}

The objectives of this project, in priority order, are:
\begin{enumerate}
    \item \textbf{Behavioral fidelity.} Produce responses that visibly follow
    MI principles (express empathy, roll with resistance, support autonomy,
    develop discrepancy) rather than defaulting to advice-giving or
    information-dumping.
    \item \textbf{State awareness.} Detect when a user shifts from openness
    to defensiveness mid-conversation and adapt the response style accordingly,
    silently -- without exposing the inferred state to the user.
    \item \textbf{Safety.} Detect crisis ideation (suicidal language,
    self-harm, overdose) and restricted-topic content (medication dosing,
    prescriptions) both before \emph{and} after model inference, and
    short-circuit to evidence-based crisis resources when warranted.
    \item \textbf{Auditability.} Produce a structured, per-turn record of every
    conversation containing the inferred state, retrieved knowledge, safety
    flags, and full text -- so that the pipeline is inspectable rather than a
    black-box LLM call.
    \item \textbf{Local deployability.} Run the production model on a CPU-only
    workstation through Ollama, with a transparent fallback to a hosted API
    when local inference is unavailable.
\end{enumerate}

\subsection{Scope and Limits}

The system is a research artifact, not a clinical product. It is not validated
under an institutional review board, has not been reviewed at scale by licensed
clinicians, and is not appropriate for unsupervised deployment to people in
acute distress. Its purpose is to demonstrate the engineering of an
MI-aligned conversational pipeline and to characterize, quantitatively, what a
small fine-tuned model can reach versus a large reference model on the same
task. The honest limits of the system are stated explicitly in
\Cref{sec:future_work}.

\subsection{Report Outline}

The remainder of this report is organized as follows.
\Cref{sec:literature} surveys related work in MI, parameter-efficient
fine-tuning, retrieval-augmented generation, and LLM-as-judge evaluation.
\Cref{sec:methodology} describes the six-stage pipeline architecture and the
fine-tuning, retrieval, and safety designs.
\Cref{sec:analysis} discusses the key engineering analyses and trade-offs --
data scarcity, model size selection, evaluation design.
\Cref{sec:experiments} specifies the evaluation set, the LLM-as-judge protocol,
and the configurations compared.
\Cref{sec:results} presents quantitative and qualitative results.
\Cref{sec:conclusions} summarizes findings, and \Cref{sec:future_work} lays out
recommended future work.
